\begin{abstract}
% The field of text-to-image generation has undergone significant advancements with the introduction of diffusion models. However, the challenge of editing real images persists, as most methods are too computationally intensive or result in poor reconstruction. In this paper, we introduce SAGE (Self-Attention Guidance for image Editing) —a novel technique that leverages pre-trained diffusion models for image editing. SAGE builds upon the DDIM algorithm and incorporates a novel guidance mechanism leveraging the self-attention layers of the diffusion U-Net. This mechanism computes a reconstruction objective by utilizing attention maps generated during the inverse DDIM process. This allows for the efficient reconstruction of unedited regions, eliminating the need for precise reconstruction of the entire input image, thereby addressing the key challenges in image editing.SAGE's excellence over other methods is not only demonstrated in quantitative and qualitative evaluations but also confirmed by a statistically-validated comprehensive user study where all 47 surveyed users preferred our approach over competing methods. Furthermore, it ranks as the top-performing method in seven out of 10 evaluated quantitative analyses, securing second and third places in the remaining three.

    The field of text-to-image generation has undergone significant advancements with the introduction of diffusion models. {\modified{}Nevertheless}, the challenge of editing real images persists, as most methods are {\modified{}either computationally intensive or produce poor reconstructions}. {\modified{}This paper introduces SAGE} (Self-Attention Guidance for image Editing) —a novel technique leveraging pre-trained diffusion models for image editing. SAGE builds upon the DDIM algorithm and incorporates a novel guidance mechanism utilizing the self-attention layers of the diffusion U-Net. This mechanism computes a reconstruction objective based on attention maps generated during the inverse DDIM process, enabling efficient reconstruction of unedited regions without the need to precisely reconstruct the entire input image. Thus, SAGE directly addresses the key challenges in image editing. The superiority of SAGE over other methods is demonstrated through quantitative and qualitative evaluations and confirmed by a statistically validated comprehensive user study, in which all 47 surveyed users preferred SAGE over competing methods. Additionally, SAGE ranks as the top-performing method in seven out of 10 quantitative analyses and secures second and third places in the remaining three.
    
\end{abstract}